Une zoologiste étudie l'évolution de la population d'une espèce animale dans un secteur géographique délimité. Elle a estimé que cette population diminue chaque année de 100 individus. Elle modélise la situation par la suite $\left(u_n\right)$ telle que $u_n$ représente le nombre d'individus au 1\up{er} mars $2025+n$.

Au 1\up{er} mars 2025, la population compte \np{2500} individus.

\begin{enumerate}
	\item
	\begin{enumerate}
		\item Donner $u_0$ à partir des données de l'exercice.
		\item Selon ce modèle, calculer le nombre d'individus au 1\up{er} mars 2026.
	\end{enumerate}

	\item
	\begin{enumerate}
		\item Exprimer $u_{n+1}$ en fonction de $u_n$ pour tout entier naturel $n$.
		\item Quelle est la nature de cette suite ? On précisera sa raison.
		\item Donner l'expression du terme général de la suite $\left(u_n\right)$ en fonction de $n$.
	\end{enumerate}

	\item Justifier le sens de variation de la suite $\left(u_n\right)$.
	\item Selon ce modèle, à partir du 1\up{er} mars de quelle année le nombre d'individus sera-t-il inférieur à \np{1250} ?
\end{enumerate}

\medskip
